\documentclass[a4j,dvipdfmx]{jsarticle}
\usepackage{amsmath,amssymb}
\usepackage{siunitx}
\usepackage{ascmac}
\usepackage{amsthm}
\usepackage{bm}
% \usepackage[dvipdfmx]{hyperref}
\begin{document}
    \part*{第一回 線型代数 問題}
    \subsection*{問題1}
        次のベクトルが一次独立か一次従属かを判定する.
        \begin{equation*}
            \begin{pmatrix}1 \\ 2 \\ 3 \end{pmatrix}, \begin{pmatrix}3 \\ 4 \\ 3 \end{pmatrix}, \begin{pmatrix}9 \\ 2 \\ 15 \end{pmatrix}
        \end{equation*}
    \subsection*{問題2}
        次の連立一次方程式をとく.
        \begin{equation*}
            (1)\left\{\begin{array}{c}
                x+2y+3z=4\\
                3x+y+4z=0\\
                -2x+3y+z=1
            \end{array}\right.\quad 
            (2)\left\{\begin{array}{c}
                3y+3z-2w=0\\
                x+y+2z+3w=0\\
                x+2y+3z+2w=0\\
                x+3y+4z+2w=0
            \end{array}\right.
        \end{equation*}
    \subsection*{問題3}
        次の行列式を計算する.
        \begin{equation*}
            (1)\begin{vmatrix}
                1 & 2 & 3 \\
                4 & 5 & 6\\
                7 & 8 & 9
            \end{vmatrix}\quad
            (2)\begin{vmatrix}
                3 & -5 & 2 & 10\\
                2 & 0 & 1 & -3\\
                -2 & 3 & 5 & 2\\
                4 & -2 & -3 & 2
            \end{vmatrix}\quad
            (3)\begin{vmatrix}
                2 & 100 & 65 & 98765421 & 15\\
                0 & 8 & 23 & 15 & 31\\
                0 & 0 & 314 & 2 & 99\\
                0 & 0 & 0 & 0.625 & -1\\
                0 & 0 & 0 & 0 &  5
            \end{vmatrix}
        \end{equation*}
    \subsection*{問題4}
        次の行列の積を計算する. ただし, $\bm{a},\bm{b},\bm{c}\in \mathbb{R}^3$
        \begin{equation*}
            (1)\begin{pmatrix}
                1 & -2 & 3\\
                2 & 1 & 0
            \end{pmatrix}
            \begin{pmatrix}
                -2 & -1 & 1\\
                1 & 2 & -3\\
                4 & 0 & 3
            \end{pmatrix}\quad 
            (2)\begin{pmatrix}
                1 & -1 & 0 & 0\\
                0 & -2 & 0 & 0\\
                0 & 0  & -2 & 3\\
                0 & 0  & 1 & 1
            \end{pmatrix}
            \begin{pmatrix}
                2 & 1 & 0 & 0\\
                0 & 1 & 0 & 0\\
                0 & 0  & 1 & 1\\
                0 & 0  & 2 & -3
            \end{pmatrix}\quad
            (3)\begin{pmatrix}
                \bm{a} & \bm{b} & \bm{c}
            \end{pmatrix}
            \begin{pmatrix}
                x\\
                y\\
                z
            \end{pmatrix}
        \end{equation*}
    \subsection*{問題5}
        次の命題を示せ.
        \begin{equation*}
            (1)A\subset A \quad (2)A\cap B\subset A
        \end{equation*}
    \subsection*{問題6}
        写像$f:\mathbb{R}\to\mathbb{R}, x\mapsto 2x+1$は全単射であることを(定義より)示せ.
\end{document}